\Needspace{8\baselineskip}
\Article{264}
In this way, then, for any given nonsquare determinant all assignable characters are divided equally into two types $P$, $Q$, so that no character in $Q$ can correspond to a properly primitive (positive) form, while for the remaining characters $P$, as far as we currently know, nothing prevents them from belonging to such forms. For these two types, note in particular the following proposition, which follows easily from their criterion: composing a character in $P$ with a character in $Q$ (by the rule of Article 246, just as if this too corresponded to a genus) gives a character in $Q$; but composing two characters in $P$. or two in $Q$ gives a character in $P$. With this theorem, half of all assignable characters can also be excluded for negative and improperly primitive genera, as follows.

I. For a negative determinant $D$, the negative genera are completely opposite to the positive ones in this respect: no character in $P$ belongs to a negative properly primitive genus, but all these genera have characters in $Q$. For when $D'\equiv1\ (\mathrm{mod.}\ 4)$, $-D'$ is a positive number of the form $4n+3$, so among $a$, $b$, $c$ etc. there are an odd number of numbers of the form $4n+3$, each of which has $-1$ as a nonresidue. Thus in this case the complete character of the form $(-1,0,D)$ contains an odd number of individual characters from $\Omega$, and so belongs to $Q$. When $D'\equiv3\ (\mathrm{mod.}\ 4)$, the same reasoning shows that among $a$, $b$, $c$ etc. there are either no numbers of the form $4n+3$, or two, or four, and so on. But since either $3,4$ or $3,8$ or $7,8$ occurs in this case among the individual characters of the form $(-1,0,D)$, its complete character also here\textperiodcentered{}belongs to $Q$. The same conclusion follows just as easily in the remaining cases
\SourceLeaf{302}
so that the negative form $(-1, 0, D)$ always has a character in $Q$. Since this form, composed with any other negative properly primitive form of the same determinant, produces such a positive form, clearly no negative properly primitive form can have a character in $P$.

II. For improperly primitive (positive) genera it is proved in the same way that the situation is either the same as for properly primitive genera or the opposite, according as $D\equiv1$ or $\equiv5\;(\mathrm{mod.}\,8)$. For in the former case we also have $D'\equiv1\;(\mathrm{mod.}\,8)$, from which it follows easily that among the numbers $a$, $b$, $c$ etc. there are either no numbers of the forms $8n+3$ and $8n+5$, or two, or four, and so on (for a product of any number of odd numbers, among which the numbers of the forms $8n+3$ and $8n+5$ together occur an odd number of times, is always either $\equiv3$ or $\equiv5\;(\mathrm{mod.}\,8)$, whereas the product of all $a$, $b$, $c$ etc. must equal either $D'$ or $-D'$). Thus the complete character of the form $(2, 1, \frac{1-D}{2})$ contains either no individual characters from $\Omega$, or two, or four, and so on, and therefore belongs to $P$. Now every improperly primitive (positive) form of determinant $D$ can be viewed as composed from $(2, 1, \frac{1-D}{2})$ and a properly primitive (positive) form of the same determinant, so clearly no improperly primitive (positive) form can have a character in $Q$ in this case. In the other case. $D\equiv5\;(\mathrm{mod.}\,8)$, everything is opposite: $D'$, which is also $\equiv5$, certainly contains an odd number of factors of the forms $8n+3$ and $8n+5$. Thus the character of the form $(2, 1, \frac{1-D}{2})$, and therefore also the character of every improperly primitive (positive) form of determinant $D$, belongs to $Q$, so no character in $P$ can correspond to an improperly primitive positive genus.

III. Finally, for a negative determinant the negative improperly primitive genera are again opposite to the positive improperly primitive genera: the former cannot have a character in $P$ or in $Q$, according as $D\equiv1$ or $\equiv5\;(\mathrm{mod.}\,8)$, or according as $-D$ is of the form $8n+7$ or $8n+3$. This follows without difficulty from the fact that composing the form $(-1, 0, D)$, whose character is in $Q$, with negative improperly primitive forms of the same determinant produces positive improperly primitive forms. Thus, when characters $Q$ are excluded from the latter, characters $P$ must be excluded from the former, and conversely.

\Needspace{11\baselineskip}
\SourceLeaf{303}
\begin{center}
\textit{A special method of decomposing primes into two squares.}
\end{center}

\Needspace{8\baselineskip}
\Article{265}
The investigations of Articles 257, 258 into the number of ambiguous classes, on which everything above is built, give many other conclusions well worth attention, which must be omitted to keep this brief. But the following is so elegant that we cannot leave it out. For a positive determinant $p$ that is a prime of the form $4n+1$, we showed that there is only one properly primitive ambiguous class. Thus all properly primitive ambiguous forms of such a determinant are properly equivalent. If, then, $b$ is the positive integer immediately below $\surd p$, and $p-bb
=a'$, the forms $(1, b, -a')$, $(-1, b, a')$ are properly equivalent; since both are clearly reduced, one occurs in the period of the other. Giving the former form index $0$ in its period, the index of the latter must be odd (since the leading coefficients of these two forms have opposite signs); set it $=2m+1$. Further, if the forms with indices $1, 2, 3$ etc. are respectively
\[
(-a', b', a''),\quad (a'', b'', -a'''),\quad (-a''', b''', a'''')\text{ etc.}
\]
then clearly the indices $2m$, $2m-1$, $2m-2$, $2m-3$ etc. correspond to the forms
\[
(a', b, -1),\quad (-a'', b', a'),\quad (a''', b'', -a''),\quad (-a'''', b''', a''')\text{ etc.}
\]
Thus, if the form of index $m$ is $(A, B, C)$, it is also $(-C, B, -A)$, so $C=-A$ and $p=BB+AA$. Consequently every prime of the form $4n+1$ can be decomposed into two squares (a proposition we derived above, in Article 182, from entirely different principles). Such a decomposition can be found by a very simple and entirely uniform method: develop the period of the reduced form whose determinant is that prime and whose leading coefficient is $1$, up to a form whose outer coefficients are equal in magnitude and opposite in sign. Thus, \textit{for example,}\ for $p=233$ we have $(1, 15, -8)$, $(-8, 9, 19)$, $(19, 10, -7)$, $(-7, 11, 16)$, $(16, 5, -13)$, $(-13, 8, 13)$, and $233=64+169$. Moreover, $A$ must clearly be odd (since $(A, B, -A)$ must be properly primitive), and so $B$ is even. --- Since for a positive determinant $p$ that is a prime of the form $4n+1$ the improperly primitive order too contains only one ambiguous class, clearly if $g$ is the
\SourceLeaf{304}
odd number immediately below $\surd p$, and $p-gg=4h$, the reduced improperly primitive forms $(2, g, -2h)$, $(-2, g, 2h)$ are properly equivalent, and therefore one occurs in the period of the other. Reasoning just as above shows that the period of the form $(2, g, -2h)$ contains a form whose outer coefficients are equal in magnitude and opposite in sign, so this too gives a decomposition of the number $p$ into two squares. The outer coefficients of this form are clearly even, and therefore its middle coefficient odd. Since a prime can be decomposed into two squares in only one way, the form found by this second method is either $(B, \pm A, -B)$ or $(-B, \pm A, B)$. Thus in our example, for $p=233$ we have $(2, 15, -4)$, $(-4, 13, 16)$, $(16, 3, -14)$, $(-14, 11, 8)$, $(8, 13, -8)$, and $233=169+64$ as above.
